%% Web edition. The GitHub edition differs only in the author line and in the Disclosure section (see github-author.patch).
\documentclass[11pt,a4paper]{article}
\usepackage{lmodern}
\usepackage[T1]{fontenc}
\usepackage[utf8]{inputenc}
\usepackage[a4paper,top=27mm,bottom=27mm,left=28mm,right=28mm]{geometry}
\usepackage{amsmath,amssymb,amsthm}
\usepackage{booktabs}
\usepackage{microtype}
\usepackage{enumitem}
\usepackage[hidelinks]{hyperref}
\usepackage{url}

\newtheorem{theorem}{Theorem}[section]
\newtheorem{lemma}[theorem]{Lemma}
\newtheorem{proposition}[theorem]{Proposition}
\newtheorem{corollary}[theorem]{Corollary}
\theoremstyle{definition}
\newtheorem{definition}[theorem]{Definition}
\theoremstyle{remark}
\newtheorem{remark}[theorem]{Remark}

\newcommand{\R}{\mathbb{R}}
\newcommand{\E}{\mathbb{E}}
\newcommand{\Var}{\operatorname{Var}}
\newcommand{\Cov}{\operatorname{Cov}}
\newcommand{\cE}{\mathcal{E}}
\newcommand{\hb}{\hbar}
\newcommand{\xc}{x_c}
\newcommand{\rcol}{\rho_{\mathrm{col}}}
\newcommand{\Csite}{C_{\mathrm{site}}}
\newcommand{\code}[1]{\texttt{#1}}
%% Distribution addresses of the Lean bundle (set by the web edition).
\newcommand{\bundleurl}{https://computoergosum.com/principia/lean/principia-ym-coarse-chain-poincare-2026-10-15.tar.gz}
\newcommand{\axiomsurl}{https://computoergosum.com/principia/lean/axioms-ym-coarse-chain-poincare-2026-10-15.txt}

\title{Convex blocks and a coarse Dobrushin chain: a Poincar\'e bound for the core of the double-well path measure with constants polynomial in $\hbar$}
\author{Yoshikazu Hosomura\\[.3em]\normalsize with the assistance of Claude (Claude Code, Anthropic)}
\date{2026-10-17 (version 2; version 1: 2026-10-15; revised 2026-10-30)}

\begin{document}
\maketitle
\begin{center}\small wvbks0 [at] gmail.com · \url{https://computoergosum.com/en/index.html}\end{center}
\begin{center}\small AI use is disclosed in \S\ref{sec:ai}; what is not claimed is in \S\ref{sec:notclaimed}.\end{center}

\begin{abstract}
We consider the Euclidean path measure of the quartic double well $V(z)=(z^2-1)^2/8$ on a time lattice of spacing $\varepsilon=1/2$, restricted to the box $[-\xc,0]^k$ with $\xc=0.89$ and with both outer boundary values in the well, $\alpha,\beta\le-\xc$; this ``core'' measure carries the weight $e^{-S/\hbar}$ and the Dirichlet form $(\hbar/\varepsilon)\sum_i|\partial_if|^2$. We prove three Poincar\'e bounds on paper. Theorem~B: a block-heat-bath path coupling and Ollivier's contraction-to-gap criterion turn a block Poincar\'e constant $C_B$ and a Dobrushin--Shlosman type average influence $\delta<1$ into $C_P\le C_B/(1-\delta)$ for every $k$. Theorem~BP: for blocks of length at most $15$ and boundary values in $(-\infty,0]$, $C_B\le 9.18+c_1\hbar^{-8}$ with $\ln c_1\le 89.38$, for all $\hbar\in(0,1]$; the exponential $e^{c/\hbar}$ is removed by a hidden convexity of the action in the variables $z_i^2$ and a one-dimensional Hardy-type bound. Theorem~CC: for the coarse chain $y_m=z_{8m}$ with Dobrushin column sum $\rcol(\hbar)<1$, $C_P\le 9.18+455.84\,\Csite(\hbar)/(1-\rcol(\hbar))$ with $\Csite\le e^{80.79}\hbar^{-8}$, uniformly in $k$. Theorems BP and CC use, in addition, the Bakry--\'Emery criterion on convex domains and the Brascamp--Lieb variance inequality, which are standard results and are cited. The hypothesis $\rcol<1$ is established for every $\hbar\in(0,1]$: on $[10^{-5},1]$ by an interval-arithmetic certificate (directed outward rounding) over $239$ intervals of $\hbar$ that cover the range without gaps, giving $\rcol\le0.400$; on $(0,10^{-5}]$ on paper, up to one one-dimensional lower bound checked by interval arithmetic, giving $\rcol\le0.874$ (Proposition~PIN). Consequently $C_P\le9.18+e^{88.99}\hbar^{-8}$ for all $\hbar\in(0,1]$, uniformly in $k$, under the stated assumptions of the certificate and two cited inequalities. The algebraic steps are written in Lean~4; the analysis is on paper. The model is one-dimensional; nothing is claimed about four-dimensional lattice gauge theory.
\end{abstract}

\section{Introduction}\label{sec:intro}

A Poincar\'e inequality for a Gibbs measure in a double-well potential typically has a constant of order $e^{\Delta/\hbar}$, with $\Delta$ the barrier height (Kramers). Restricting to a single well removes the barrier, but the restricted measure is not log-concave: the potential $V$ has $V''<0$ on $(-1/\sqrt3,0]$ and the box reaches up to the top of the barrier at $0$. The standard routes from a Dobrushin--Shlosman condition to a Poincar\'e or log-Sobolev inequality bound block constants by the Holley--Stroock perturbation lemma [HS] and therefore pay a factor exponential in the interaction strength $1/(\varepsilon\hbar)$ (for instance the block constant $c_0e^{4|X|\|\Phi\|}$ in [SZ]; see also [Ma]). The aim of this paper is a Poincar\'e constant for the single-well core that is \emph{polynomial} in $\hbar$ and uniform in the length $k$.

The question arises in an online textbook by the authors [KSR], where a Poincar\'e bound of this form for the core (the \emph{core Poincar\'e bound} below) is one of the inputs on which a conditional argument for the double-well path measure rests. A companion note [KSR2] proves the opposite direction: an inverse squared spectral gap bounds the Langevin Poincar\'e constant of Euclidean path measures \emph{from below}. The present paper concerns the \emph{upper} bound for the core, which is a different object; the two results are not used in each other's proofs.

\paragraph{What the paper does.}
\begin{itemize}[leftmargin=*]
\item Theorem~\ref{thm:B} (\S\ref{sec:B}) is an abstract step: a path coupling for the block heat bath, combined with Ollivier's criterion [Ol, Prop.~29, Cor.~30], gives $C_P\le C_B/(1-\delta)$. No step introduces $e^{1/(\varepsilon\hbar)}$.
\item Theorem~\ref{thm:BP} (\S\ref{sec:BP}) bounds the block constant polynomially in $\hbar$.
\item Theorem~\ref{thm:CC} (\S\ref{sec:CC}) replaces the block chain by a one-site Dobrushin chain on every eighth site. It needs neither Theorem~\ref{thm:B} nor Theorem~\ref{thm:BP} as stated, only their ingredients, and its single numerical hypothesis is $\rcol(\hbar)<1$.
\item Proposition~\ref{prop:DOB} (\S\ref{sec:DOB}) establishes $\rcol(\hbar)<1$ for every $\hbar\in(0,1]$, with the strength of each part named in the statement: interval arithmetic on $[10^{-5},1]$, paper on $(0,10^{-5}]$ (Proposition~\ref{prop:PIN}).
\end{itemize}

\paragraph{Changes from the version of 2026-10-15.} That version certified $\rcol<1$ at $23$ values of $\hbar$ by a double-precision computation with outward margins (not interval arithmetic) and left the values between the points, and $\hbar<10^{-5}$, as a numerical reading. The present version replaces the point certificate by an interval-arithmetic certificate over $239$ intervals covering $[10^{-5},1]$ (\S\ref{sec:discipline}), adds Proposition~\ref{prop:PIN} for $\hbar\le10^{-5}$, and extends the conclusion to all of $(0,1]$ (Corollary~\ref{cor:points}). Two Lean modules were added (\S\ref{sec:lean}).

\section{Setting}\label{sec:setting}

Fix $\varepsilon=1/2$, $V(z)=(z^2-1)^2/8$, $\xc=0.89$ and $\hbar\in(0,1]$. For $k\ge1$ and outer boundary values $\alpha,\beta\in\R$ put $z_0=\alpha$, $z_{k+1}=\beta$ and
\[
S(z)=\sum_{i=0}^{k}\frac{(z_{i+1}-z_i)^2}{2\varepsilon}+\sum_{i=1}^{k}\varepsilon V(z_i),\qquad
\mu_k^{\alpha,\beta,\hbar}(dz)=\frac{1}{Z}\,e^{-S(z)/\hbar}\,\mathbf 1_{\Omega_k}(z)\,dz,\quad \Omega_k=[-\xc,0]^k .
\]
The Dirichlet form is $\cE(f)=(\hbar/\varepsilon)\sum_{i=1}^k\mu|\partial_if|^2$ (Neumann conditions on $\partial\Omega_k$), and $C_P(\mu)$ is the least constant with $\Var_\mu f\le C_P(\mu)\,\cE(f)$ for Lipschitz $f$. The \emph{core} is the family $Q_k=\mu_k^{\alpha,\beta,\hbar}$ with
\[
\textbf{(END)}\qquad \alpha,\beta\le-\xc ,
\]
that is, both outer neighbours sit in the well. (For $a>0$ the boundary term is not compatible with the hidden convexity below; the restriction (END) is part of every statement.)

The Hessian of $S$ in the $z$ coordinates is $\varepsilon^{-1}T_n+\varepsilon\,\mathrm{diag}\,V''(z_i)$ with $T_n=\mathrm{tridiag}(-1,2,-1)$, and $V''\ge V''(0)=-1/2$. Hence on any block of $n$ consecutive free sites, for arbitrary real boundary values,
\begin{equation}\label{eq:kappa}
\mathrm{Hess}\,S\ \ge\ \kappa(n)\,I,\qquad \kappa(n):=\tfrac{15}{4}-4\cos\!\big(\tfrac{\pi}{n+1}\big).
\end{equation}
Numerically $\kappa(6)=0.1461$, $\kappa_0:=\kappa(7)=0.054482$ and $\kappa(8)=-0.00877$. Blocks of length $\le7$ are uniformly convex in $z$; length $8$ is the first that is not. Every construction below cuts the chain into pieces of length $\le7$.

\paragraph{Hidden convexity (K$'$).} For $z\le0$ put $\rho_i=z_i^2$. Then $(z_i-z_j)^2=\rho_i+\rho_j-2\sqrt{\rho_i\rho_j}$ is jointly convex in $(\rho_i,\rho_j)$ (the geometric mean is concave) and $\varepsilon V(-\sqrt\rho)=\varepsilon(\rho-1)^2/8$ is convex. Hence, for boundary values $a,b\le0$, $\tilde S(\rho):=S(-\sqrt\rho)$ is convex on $[0,\xc^2]^n$ (indeed $(\varepsilon/4)$-strongly convex), jointly in $(\rho,a^2,b^2)$, and its minimiser is unique (Lemma~\ref{lem:K}, Appendix~\ref{app:K}). Its algebraic core is in Lean (\S\ref{sec:lean}). It is used below only through partial minimisation: if $\tilde S$ is jointly convex then $\sigma\mapsto\min_{\rho'}\tilde S(\rho',\sigma)$ is convex.

\paragraph{The coarse chain.} For $k\ge8$ let $M=\lfloor k/8\rfloor$, $y_m:=z_{8m}$ ($1\le m\le M$), and let $B_0=\{1,\dots,7\}$, $B_m=\{8m+1,\dots,\min(8m+7,k)\}$ be the blocks between coarse sites. With the coarse sites frozen, the blocks are independent (Markov property) and each is $\kappa_0$-convex by \eqref{eq:kappa}. The law $\nu$ of $y$ is a nearest-neighbour chain on $[-\xc,0]^M$ with effective pair interaction $w_B(a,b)=-\hbar\ln\int e^{-S_B(z;a,b)/\hbar}dz$. The conditional law $\nu(dy_m\mid y_{m-1}=a,\,y_{m+1}=b)$ is the marginal of the middle site $z_8$ of a chain $z_1,\dots,z_N$, $N=8+r$, with boundary values $z_0=a\in[-\xc,0]$ and $z_{N+1}=b$; in the interior $r=7$ and $b\in[-\xc,0]$, at the right end of the chain $r\in\{0,\dots,7\}$ and $b=\beta\le-\xc$. Its \emph{influence} (a Dobrushin coefficient [Do]) is
\begin{gather*}
\rho(a,b;r):=\partial_a\E[z_8\mid a,b]=\frac{\Cov(z_1,z_8)}{\varepsilon\hbar},\\
\rcol(\hbar):=\sup_m\sum_{m'=m\pm1}\sup_{a,b}\rho_{m'\leftarrow m}\ \le\ 2\sup_{r,a,b}\rho(a,b;r).
\end{gather*}

\paragraph{What is shown and what is not.} Theorems~\ref{thm:B}, \ref{thm:BP} and \ref{thm:CC} are proved on paper, with cited ingredients listed in Table~\ref{tab:status}. Theorem~\ref{thm:CC} is conditional on $\rcol(\hbar)<1$. That hypothesis is established for every $\hbar\in(0,1]$ (Proposition~\ref{prop:DOB}), under the assumptions (A1)--(A5) of the interval-arithmetic certificate (\S\ref{sec:discipline}) on $[10^{-5},1]$ and on paper on $(0,10^{-5}]$. Accordingly, for the core in which both outer values lie in the well (the compartment-side core), the core Poincar\'e bound moves from an assumption to a statement that is conditional on the assumptions of the certificate and on two cited inequalities (Table~\ref{tab:status}); it does not become a paper theorem. The box-side core of that argument, in which an outer neighbour lies beyond the summit and the conditional law is bistable with constants of type $e^{\Delta S/\hbar}$, is not covered by the hypothesis $\alpha,\beta\le-\xc$ and is not addressed here. The four-dimensional lattice gauge problem that motivates the question is not addressed (\S\ref{sec:notclaimed}).

\section{Theorem B: block heat bath and contraction}\label{sec:B}

Fix a block length $L$ and put $B_j=[j,j+L-1]\cap[1,k]$ for $j=2-L,\dots,k$, so that there are $N=k+L-1$ blocks, truncated blocks included, and every site lies in exactly $L$ of them. Let $E_B$ be the conditional expectation given the coordinates outside $B$, and $P=N^{-1}\sum_jE_{B_j}$ the block heat bath. For a site $y$ and a block $B$ having $y$ as a boundary site, let $J_B(y)=\sup_s\sum_{i\in B}\partial_s\E_B^{s,\cdot}[z_i]$, the supremum running also over the other boundary value of $B$, and
\[
\delta:=\sup_y\frac{J_{\mathrm{right}}(y)+J_{\mathrm{left}}(y)}{L}.
\]

\begin{theorem}[Theorem B]\label{thm:B}
Let $\alpha,\beta\le-\xc$ and $\hbar>0$. Assume
\textup{(BP)} every block $B_j$ \textup{(}length $\le L$, all boundary values in $(-\infty,0]$\textup{)} satisfies $\Var_{E_B}(f)\le C_B(\hbar/\varepsilon)\sum_{i\in B}E_B|\partial_if|^2$, and
\textup{(DS)} $\delta<1$.
Then for every $k\ge1$,
\[
C_P\big(\mu_k^{\alpha,\beta,\hbar}\big)\le\frac{C_B}{1-\delta}.
\]
The proof uses the FKG lattice condition and Holley's criterion for monotone coupling [Ho], [Pr] \textup{(}standard results, cited\textup{)} and Ollivier's criterion [Ol] \textup{(}checked against the arXiv version\textup{)}.
\end{theorem}

\begin{proof}
\emph{Monotone coupling.} The cross term of $-S/\hbar$ is $+z_iz_{i+1}/(\varepsilon\hbar)$, supermodular, and the reference measure is a product of intervals; the block measures satisfy the FKG lattice condition and the density ratio between boundary values $s<s'$ is $\exp((s'-s)z_{\mathrm{first}}/(\varepsilon\hbar))$, increasing. By Holley's criterion and Strassen's theorem there is a coupling with $z\le z'$ coordinatewise, and then $W_1(\mu_B^{s},\mu_B^{s'})=\sum_i(\E^{s'}z_i-\E^{s}z_i)\le|s'-s|\,J_B(y)$.

\emph{Path coupling.} If $z,z'$ differ only at the site $y$, by $d$, then one step of $P$ chooses a block containing $y$ ($L$ choices; the two configurations are coupled to agree), a block having $y$ as a boundary site (two choices; the discrepancy is at most $d+J_B(y)d$), or another block (discrepancy $d$). Hence $W_1(P(z,\cdot),P(z',\cdot))\le(1-\kappa)d$ with $\kappa=(L/N)(1-\delta)$; for general $z,z'$ change one coordinate at a time (the box is a product) and use the triangle inequality.

\emph{Contraction to gap.} $P$ is self-adjoint and positive in $L^2(\mu)$. Since $\Omega_k$ is compact, $\|P^nf-\mu f\|_2\le\mathrm{diam}(\Omega_k)(1-\kappa)^n\|f\|_{\mathrm{Lip}}$ for Lipschitz $f$, so the spectral measure of $f-\mu f$ has no mass in $(1-\kappa,1]$; by density of Lipschitz functions, $\langle f,(I-P)f\rangle\ge\kappa\Var f$ for all $f$ [Ol, Prop.~29, Cor.~30].

\emph{Comparison.} $\langle f,(I-P)f\rangle=N^{-1}\sum_j\mu(\Var_{B_j}f)\le(C_BL/N)\,\cE(f)$ by (BP), since each site lies in $L$ blocks. Combining, $\Var f\le C_B(1-\delta)^{-1}\cE(f)$.
\end{proof}

\begin{remark}
The ingredients are of known types (path coupling under a Dobrushin--Shlosman type condition [Ma], comparison of block dynamics, Ollivier's criterion). We do not know whether this combination appears in the literature in this form; in the references we were able to read, the block constant is controlled by a Holley--Stroock perturbation, which is where $e^{c/(\varepsilon\hbar)}$ enters. Here that step is the hypothesis (BP), and Theorem~\ref{thm:BP} supplies it. Truncated blocks enter the supremum in $\delta$.
\end{remark}

The algebraic steps (the one-step contraction inequality, the passage from spectral weight to gap, and the final comparison) are in Lean (\S\ref{sec:lean}). The quantity $\delta$ itself is not bounded on paper in this paper; Theorem~\ref{thm:CC} below avoids it.

\section{Theorem BP: block constants polynomial in \texorpdfstring{$\hbar$}{hbar}}\label{sec:BP}

For a block of length $L\le15$ with boundary values $a,b\in(-\infty,0]$ write $\mu_{a,b}\propto e^{-S/\hbar}\mathbf 1_{[-\xc,0]^L}$ and $C_B(\hbar):=\sup_{L\le15,\,a,b\le0}C_P(\mu_{a,b})$.

\begin{lemma}[short blocks; BE$_7$]\label{lem:BE7}
If $n\le7$, then for every $\hbar>0$ and all real boundary values, $C_P(\mu_{a,b})\le\varepsilon/\kappa(n)\le\varepsilon/\kappa_0=9.1774$.
\end{lemma}
\begin{proof}
By \eqref{eq:kappa} the potential $S/\hbar$ is $(\kappa(n)/\hbar)$-strongly convex on the convex set $[-\xc,0]^n$. The Bakry--\'Emery criterion for a strongly log-concave measure restricted to a convex set (Neumann boundary) gives $\Var f\le(\hbar/\kappa(n))\int|\nabla f|^2d\mu$ [BGL], [BL, Thm.~4.1] \textup{(}the criterion is standard and is cited; its application to the box $[-\xc,0]^n$ through \eqref{eq:kappa} is the step of this paper\textup{)}.
\end{proof}

\begin{lemma}[variance split at the middle site]\label{lem:split}
Let $8\le L\le15$, $j=\lceil L/2\rceil$, let $m$ be the law of $z_j$ and $C_m$ its Poincar\'e constant for $\Var_mg\le C_m\int|g'|^2dm$. Then $C_P(\mu_{a,b})\le\varepsilon/\kappa_0+\big(4/(\varepsilon\kappa_0^2)\big)\,C_m/\hbar$, with $4/(\varepsilon\kappa_0^2)=2695.2$.
\end{lemma}
\begin{proof}
Both sides of $j$ have length $\le7$. By the law of total variance, $\Var_\mu f=m(\Var_{\mu_s}f)+\Var_m g$, $g(s)=\mu_s(f)$, where $\mu_s$ is the law with $z_j=s$ frozen, a product of two $\kappa_0$-convex blocks; Lemma~\ref{lem:BE7} and tensorisation give the first term $\le(\hbar/\kappa_0)A$, $A=\sum_{i\ne j}\mu|\partial_if|^2$. Differentiating, $g'(s)=\mu_s(\partial_jf)+\Cov_{\mu_s}(f,h)/(\varepsilon\hbar)$ with $h=z_{j-1}+z_{j+1}$; Cauchy--Schwarz and $\Var_{\mu_s}h\le2\hbar/\kappa_0$ give $|g'|^2\le2\mu_s|\partial_jf|^2+4A_s/(\varepsilon^2\kappa_0^2)$. Integrate against $m$ and rescale to $\cE$.
\end{proof}

\begin{lemma}[one-dimensional Hardy-type bound; H1D]\label{lem:H1D}
Let $m(dt)=Z^{-1}e^{-F(t)/\hbar}dt$ on $I=[-\ell,0]$ with $F$ continuous, and $\Delta:=\sup_{y\le s\le x}\big[F(s)-\max(F(x),F(y))\big]$. Then $C_m\le2\ell^2e^{\Delta/\hbar}$.
\end{lemma}
\begin{proof}
$\Var_mg=\iint_{y<x}(g(x)-g(y))^2dm\,dm$ and $(g(x)-g(y))^2\le\int_y^x|g'|^2p\int_y^xp^{-1}$. By Fubini $\Var_mg\le\int|g'(t)|^2p(t)K(t)dt$ with $K(t)=\iint_{y\le t\le x}\int_y^xp(x)p(y)/p(s)\,ds\,dx\,dy$. For $y\le s\le x$, $p(x)p(y)/p(s)\le e^{\Delta/\hbar}(p(x)+p(y))$, and $\iint_{y\le t\le x}(x-y)(p(x)+p(y))\le2\ell^2$.
\end{proof}
If $F$ is quasi-convex ($\Delta=0$), $C_m\le2\ell^2$ independently of $\hbar$: the factor $e^{c/\hbar}$ can only come from a second valley of the free energy. The discrete form of Lemma~\ref{lem:H1D} is in Lean; the Fubini step is on paper. We do not use Muckenhoupt's criterion [Mu].

\begin{lemma}[quasi-convex free energy plus small oscillation; QC]\label{lem:QC}
Let $n=L-1\le14$, $F_\hbar(s)=-\hbar\log\int e^{-S(z',s)/\hbar}dz'$ the free energy of the middle site, $F_0(s)=\min_{z'}S(z',s)$ and $R=F_\hbar-F_0$. Then \textup{(a)} $F_0$ is quasi-convex on $[-\xc,0]$, and \textup{(b)} $R(s)-R(s')\le\hbar\log(K_1\hbar^{-n/2})$ for all $s,s'$, where $K_1=(2\pi/\kappa_0)^{n/2}G^ne^{n(1+\Lambda/(2G^2))}$, $G=6\xc/\varepsilon+\varepsilon\max|V'|=10.776$, $\Lambda=4/\varepsilon+\varepsilon\max V''=8.344$. Consequently $\Delta(F_\hbar)\le\hbar\log(K_1\hbar^{-n/2})$.
\end{lemma}
\begin{proof}
(a) $F_0(s)=\tilde F_0(s^2)$ with $\tilde F_0(\sigma)=\min_{\rho'}\tilde S(\rho',\sigma)$ convex by (K$'$) and partial minimisation; $s\mapsto s^2$ is monotone on $[-\xc,0]$, so the sublevel sets of $F_0$ are intervals.
(b) Let $z^{(s)}$ be the minimiser (unique by Lemma~\ref{lem:BE7}), $g=\nabla_{z'}S(z^{(s)},s)$ and $d=z'-z^{(s)}$. The KKT condition gives $\langle g,d\rangle\ge0$ on the box, and $(\kappa_0/2)|d|^2+\langle g,d\rangle\le S-F_0\le(\Lambda/2)|d|^2+\langle g,d\rangle$. The upper bound gives $\int e^{-(S-F_0)/\hbar}\le(2\pi\hbar/\kappa_0)^{n/2}$; on the cube obtained by moving each coordinate inward by $\hbar/G$ one has $\langle g,d\rangle\le n\hbar$, whence $\int e^{-(S-F_0)/\hbar}\ge(\hbar/G)^ne^{-n(1+\Lambda\hbar/(2G^2))}$ when $|a|,|b|\le3\xc$; for $|a|>3\xc$ the first coordinate of the minimiser is pinned at $-\xc$ and the same estimate holds with the first direction rescaled by $\hbar/g_1$. A quasi-convex function plus an oscillation of size $D$ has barrier at most $D$.
\end{proof}

\begin{theorem}[Theorem BP]\label{thm:BP}
For all $\hbar\in(0,1]$, $L\le15$ and $a,b\in(-\infty,0]$,
\[
C_P(\mu_{a,b})\le\frac{\varepsilon}{\kappa_0}+c_1\hbar^{-8},\qquad \frac{\varepsilon}{\kappa_0}=9.1774,\quad \ln c_1\le89.38 ,
\]
and $C_P(\mu_{a,b})\le\varepsilon/\kappa(L)\le9.18$ if $L\le7$. The bound uses \textup{(K$'$)} and Lemma~\ref{lem:BE7}, whose convex-domain Bakry--\'Emery step is cited.
\end{theorem}
\begin{proof}
Lemmas~\ref{lem:H1D} and \ref{lem:QC} give $C_m\le2\xc^2K_1\hbar^{-n/2}$, and Lemma~\ref{lem:split} gives $C_P\le\varepsilon/\kappa_0+(8\xc^2/(\varepsilon\kappa_0^2))K_1\hbar^{-1-n/2}$. For $n\le14$ and $\hbar\le1$, $\hbar^{-1-n/2}\le\hbar^{-8}$ and $K_1$ is largest at $n=14$: $\ln(8\xc^2/(\varepsilon\kappa_0^2))=8.359$ and $\ln K_1=33.234+33.283+14.503=81.020$, so $\ln c_1\le89.379$.
\end{proof}

The constant is crude. (K$'$) is used once, in Lemma~\ref{lem:QC}(a), to exclude a second valley of $F_\hbar$; without it $\Delta=O(1)$ would be possible and the bound would be of Kramers type $e^{\Delta/\hbar}$. With Theorem~\ref{thm:B}, Theorem~\ref{thm:BP} gives $C_P(Q_k)\le(9.18+c_1\hbar^{-8})/(1-\delta)$ whenever the averaged influence $\delta$ for $L=15$ is $<1$; we do not certify $\delta$ here and use Theorem~\ref{thm:CC} instead.

\section{Theorem CC: the coarse Dobrushin chain}\label{sec:CC}

Let $M_7=\varepsilon^{-1}T_7-(\varepsilon/2)I$ (the worst-case Hessian of a block of length $7$), so $\lambda_{\min}(M_7)=\kappa_0$, and
\[
g_7:=(M_7^{-1})_{11}=(M_7^{-1})_{77}=\tfrac{1832324}{1770735}=1.03478 ,
\]
computed by the continued fraction of the tridiagonal Schur recursion (value in Lean). For a block $B$ of length $n\le7$, $(M_n^{-1})_{11}\le g_7$.

\begin{lemma}[stochastic monotonicity of the coarse conditional laws]\label{lem:B1}
$\nu(\cdot\mid a,b)$ is increasing in $a$ in the likelihood-ratio order, so that $W_1(\nu(\cdot\mid a,b),\nu(\cdot\mid a',b))=|\E_{a'}y_m-\E_ay_m|\le\sup\rho\cdot|a'-a|$.
\end{lemma}
\begin{proof}
The law of $z_8$ given $a$ is the image of the Markov chain of Lemma~\ref{lem:R} below, whose kernels are increasing in their starting point in the likelihood-ratio order (Lemma~\ref{lem:M}(a)); composition of TP$_2$ kernels is TP$_2$ [Ka]. Hence the law of $z_8$ is increasing in $a$, and the monotone coupling realises $W_1$ as the difference of means.
\end{proof}

\begin{theorem}[coarse chain]\label{thm:C}
If $\rcol<1$, then for every Lipschitz $g$ on $[-\xc,0]^M$,
\[
\Var_\nu g\le\frac{\Csite}{1-\rcol}\,\frac{\hbar}{\varepsilon}\sum_m\nu|\partial_{y_m}g|^2,\qquad \Csite:=\sup_{m,a,b}\frac{\varepsilon}{\hbar}C_P(\nu(\cdot\mid a,b)).
\]
\end{theorem}
\begin{proof}
This is Theorem~\ref{thm:B} with blocks of one coarse site. Let $P=M^{-1}\sum_mP_m$ be the one-site heat bath. If $y,y'$ differ only at $m$, by $d$, choosing $m$ couples them to agree, choosing $m'=m\pm1$ costs at most $d+\rho_{m'\leftarrow m}d$ by Lemma~\ref{lem:B1}, and other choices leave $d$; so $W_1(P(y,\cdot),P(y',\cdot))\le(1-(1-\rcol)/M)d$. Ollivier's criterion gives $\langle g,(I-P)g\rangle\ge((1-\rcol)/M)\Var_\nu g$, and $\langle g,(I-P)g\rangle=M^{-1}\sum_m\nu(\Var_mg)\le M^{-1}\Csite(\hbar/\varepsilon)\sum_m\nu|\partial_mg|^2$.
\end{proof}

The conditional law $\nu(\cdot\mid a,b)$ is the marginal of the middle site of a chain of at most $15$ sites whose two halves have length $\le7$. Lemmas~\ref{lem:H1D} and \ref{lem:QC} (with (K$'$) extended to the boundary values, \S\ref{sec:setting}) therefore give
\begin{equation}\label{eq:Csite}
\Csite\le2\varepsilon\xc^2K_1(14)\,\hbar^{-8}=e^{80.79}\hbar^{-8}.
\end{equation}

\begin{theorem}[Theorem CC]\label{thm:CC}
Let $\varepsilon=1/2$, $V(z)=(z^2-1)^2/8$, $\xc=0.89$, $\hbar\in(0,1]$ and $\alpha,\beta\le-\xc$. Assume
\textup{(DOB)} $\rcol(\hbar)<1$.
Assume moreover the cited statements \textup{BE}$_7$ \textup{(}Lemma~\ref{lem:BE7}\textup{)} and \textup{BL}: for a measure on a convex set with $\mathrm{Hess}(S/\hbar)\ge M/\hbar$, $\Var\langle e,z\rangle\le\hbar\langle e,M^{-1}e\rangle$ [BL, Thm.~4.1] \textup{(}a standard result, cited\textup{)}. Then for every $k\ge1$,
\[
C_P(Q_k)\le\frac{\varepsilon}{\kappa_0}+c_6\,\frac{\Csite(\hbar)}{1-\rcol(\hbar)}\le 9.18+\frac{455.84\,e^{80.79}\,\hbar^{-8}}{1-\rcol(\hbar)},\qquad c_6:=\max\Big(3,\frac{6g_7}{\varepsilon^2\kappa_0}\Big)=455.84 .
\]
In particular, if $\rcol(\hbar)\le1/2$ then $C_P(Q_k)\le9.18+e^{87.61}\hbar^{-8}$, uniformly in $k$.
\end{theorem}

\begin{proof}
For $k\le7$ the core is one block and Lemma~\ref{lem:BE7} gives $C_P\le9.18$. Let $k\ge8$, write $A=\sum_{i\ \mathrm{not\ coarse}}\mu|\partial_if|^2$ and $B=\sum_m\mu|\partial_{8m}f|^2$, so $\cE(f)=(\hbar/\varepsilon)(A+B)$.
\begin{enumerate}[leftmargin=*,label=\arabic*.]
\item $\Var f=\nu(\Var(f\mid y))+\Var_\nu g$, $g(y)=\E[f\mid y]$.
\item Given $y$ the blocks are independent and $\kappa_0$-convex; by Lemma~\ref{lem:BE7} and tensorisation, $\nu(\Var(f\mid y))\le(\varepsilon/\kappa_0)(\hbar/\varepsilon)A$.
\item $\partial_{y_m}g=\E[\partial_{8m}f\mid y]+[\Cov(f,z_{8m-1}\mid y)+\Cov(f,z_{8m+1}\mid y)]/(\varepsilon\hbar)$, since $\partial_{y_m}S=(2y_m-z_{8m-1}-z_{8m+1})/\varepsilon+\varepsilon V'(y_m)$ and terms depending on $y_m$ alone do not contribute to covariances.
\item $X=z_{8m\pm1}$ depends on one block $B$ only, so $\Cov(f,X\mid y)=\Cov(f_B,X\mid y)$ with $f_B=\E[f\mid y,x_B]$. Then $\Var(f_B\mid y)\le(\hbar/\kappa_0)A_B$ (block Poincar\'e and Jensen) and $\Var(X\mid y)\le\hbar g_7$ (BL), so $(\Cov/(\varepsilon\hbar))^2\le g_7A_B/(\varepsilon^2\kappa_0)$.
\item $(u+v+w)^2\le3(u^2+v^2+w^2)$ and the fact that each block touches at most two coarse sites give $\sum_m\nu|\partial_mg|^2\le3B+(6g_7/(\varepsilon^2\kappa_0))A\le c_6(A+B)$.
\item Theorem~\ref{thm:C}: $\Var_\nu g\le(\Csite/(1-\rcol))\,c_6\,(\hbar/\varepsilon)(A+B)$.
\item Adding 2.\ and 6.\ gives the claim; \eqref{eq:Csite} gives the numerical form, and $\ln(c_6e^{80.79})=86.9121$, so $86.9121+\ln2<87.61$.\qedhere
\end{enumerate}
\end{proof}

Theorem~\ref{thm:CC} uses neither Theorem~\ref{thm:B} with $L=15$ nor Theorem~\ref{thm:BP} as stated; it uses their ingredients (path coupling, Lemmas~\ref{lem:BE7}, \ref{lem:H1D}, \ref{lem:QC}). The factor $\rcol$ enters only through $1/(1-\rcol)$; its numerical role is that of a single hypothesis, examined next.

\section{The Dobrushin certificate}\label{sec:DOB}

\subsection{Why a certificate}
A paper bound for $\rho$ via a sup of $\partial_s\E[z_1\mid z_8=s]$ times $\Var(z_8)/(\varepsilon\hbar)$ does not close at small $\hbar$, and the curvature route cannot close it in principle:

\begin{proposition}[curvature floor]\label{prop:N1}
For all boundary values with both neighbours of the middle site in the box and all $s\in[-\xc,0]$, $F_\hbar''(s)\le2/\varepsilon+\varepsilon\max V''=4.3441$. Hence every bound of the form $\Var_m\le\hbar/\kappa$ with $\kappa\le\inf F_\hbar''$ gives $\Var_m/(\varepsilon\hbar)\ge0.4604$.
\end{proposition}
\begin{proof}
$F_\hbar''=\E[\partial_s^2S\mid s]-\Var(\partial_sS\mid s)/\hbar\le2/\varepsilon+\varepsilon V''(s)$, and $\max_{[-\xc,0]}V''=V''(-\xc)=0.68815$.
\end{proof}
At small $\hbar$ the middle site is pressed against the wall $-\xc$ by the slope $\varepsilon V'(-\xc)=0.0463$ and its variance is of order $\hbar^2$, a boundary effect that curvature bounds do not see. The certificate below follows the law of the chain step by step instead.

\subsection{Two lemmas}
Let $\lambda=1/(\varepsilon\hbar)$, $K(x)=e^{-\lambda x^2/2}$ and $[lo,hi]=[-\xc,0]$. Define backward functions $g_N(v)=e^{-\varepsilon V(v)/\hbar}K(b-v)$, $g_i(v)=e^{-\varepsilon V(v)/\hbar}\int K(w-v)g_{i+1}(w)\,dw$.

\begin{lemma}[influence recursion; Lemma R]\label{lem:R}
Given $z_0=a$, the chain $(z_1,\dots,z_N)$ is Markov with kernels $P_i(u,dv)=K(v-u)g_i(v)dv/\!\int K(v'-u)g_i(v')dv'$ \textup{(}Doob $h$-transform\textup{)}. Put $f_8=\mathrm{id}$, $f_{i-1}=P_if_i$; then $\E[z_8\mid a]=f_0(a)$, $\rho=f_0'(a)$, and
\[
f_{i-1}'(u)=\lambda\int f_i'(t)\,k_i(t;u)\,dt,\qquad k(t):=\Cov(V,\mathbf 1\{V>t\})=\bar F(t)\!\int_{lo}^{t}\!F+F(t)\!\int_t^{hi}\!\bar F\ \ge0 ,
\]
where $V\sim P_i(u,\cdot)$ with distribution function $F$ and $\bar F=1-F$.
\end{lemma}
\begin{proof}
$\partial_u(P_if)(u)=\lambda\Cov_{P_i(u,\cdot)}(V,f(V))$ because $\partial_u\log K(v-u)=\lambda(v-u)$; and $f(V)=f(lo)+\int f'(t)\mathbf 1\{V>t\}dt$ gives $\Cov(V,f(V))=\int f'(t)k(t)dt$ (Hoeffding [Hoe]).
\end{proof}
By induction $f_i'\ge0$, without appeal to FKG. If $f_i'\le\Phi_i$ then $f_{i-1}'\le\lambda\int\Phi_ik_i$, and an upper bound for $k$ follows from two-sided bounds $F_{lo}\le F\le F_{hi}$ alone, as a sum of products of nonnegative terms with no cancellation. Explicitly, for $t\in[t_j,t_{j+1}]$,
\[
k(t)\le\big(1-F_{lo}(t_j)\big)\int_{lo}^{t_{j+1}}F_{hi}+F_{hi}(t_{j+1})\int_{t_j}^{hi}\big(1-F_{lo}\big),
\]
and the lower bound has the same form with the roles of $F_{lo}$ and $F_{hi}$ exchanged. For $V\equiv0$ and an unbounded box the product of the eight factors is $\prod_{i=1}^8(16-i)/(17-i)=1/2$, the exact value for the Gaussian bridge.

\begin{lemma}[sandwiching; Lemma M]\label{lem:M}
\textup{(a)} $P_i(u,\cdot)$ is increasing in $u$ in the likelihood-ratio order. \textup{(b)} $g_i^{(b)}(v)$ is TP$_2$ in $(v,b)$, so $P_i^{(b)}(u,\cdot)$ is increasing in $b$; the limit $b\to-\infty$ is the chain with $z_N=lo$. \textup{(c)} Write $P_i(u,dw)\propto\exp(\lambda uw-\lambda w^2+\chi_i(w))$. If on each cell of a partition $\chi_i'\in[s_j^-,s_j^+]$, then replacing $\chi_i$ by the piecewise linear functions with slopes $s_j^\pm$ \textup{(}jumping only upwards, resp.\ downwards, at nodes\textup{)} gives laws $L^-\le P_i(u,\cdot)\le L^+$ in the likelihood-ratio order.
\end{lemma}
\begin{proof}
(a) The density is $e^{\lambda uv}$ times a function of $v$. (b) $g_N$ is $e^{\lambda bv}$ times functions of $v$ and of $b$; $K(w-v)$ is TP$_2$; compositions of TP$_2$ kernels are TP$_2$ by the $2\times2$ Cauchy--Binet formula [Ka]. (c) The log-ratio of the densities is monotone on each cell and across nodes.
\end{proof}
The Gaussian factor $e^{\lambda uw-\lambda w^2}$ times an exponential of a linear function is integrated exactly on each cell by the error function, so the strong push towards the wall at small $\hbar$ costs nothing in cell width; the remaining coarseness is the slope range $s_j^+-s_j^-$.

\paragraph{The certificate inequality.} On a cell $[u_l,u_{l+1}]$ of $a$ and a cell $[b_{lo},b_{hi}]$ of $b$ one has $F_{L^+_i(u_{l+1})}\le F_{P_i}\le F_{L^-_i(u_l)}$, with cell masses given by closed erf formulas. Inserting these into Lemma~\ref{lem:R} and running $\Phi_i$ backward as cellwise constants yields
\[
\sup_{a\in[-\xc,0],\,b\in[b_{lo},b_{hi}]}\rho(a,b)\le\max_l\Phi_0(l).
\]
The variable $a$ is covered by cells continuously and $b$ by the endpoint slopes of its cells; no reading between grid points of $a$ or $b$ remains. Reversing the roles gives a lower bound, used only as a check.

\subsection{The certificate in interval arithmetic}\label{sec:discipline}
Every quantity that enters an inequality of the certificate is enclosed in an interval with outward rounding. The implementation (Python, double precision) rests on the following assumptions, which are part of the statement of Proposition~\ref{prop:DOB}(a).
\begin{enumerate}[leftmargin=*,label=(A\arabic*)]
\item IEEE~754 double precision with round-to-nearest; after each of the five basic operations the result $r$ is widened to $r\pm(|r|2^{-51}+2^{-1074})$, which contains the exact result.
\item Constants and the coefficient tables are computed with the interval type of \code{mpmath} [mp] and stored as enclosures.
\item $\mathrm{erfcx}$ on $[0,16)$ is evaluated from degree-$8$ Taylor tables at the nodes $k/32$ with an enclosure of the remainder (the derivatives are completely monotone); for $x\ge16$ from the asymptotic series, whose remainder is bounded by the first omitted term [DLMF, \S7.12(i)]; $\mathrm{erf}$ for $|x|\le0.3$ from the alternating series, and otherwise from $\mathrm{erfcx}$.
\item Polynomials and sums are evaluated in floating point once, with an a priori bound on the rounding error of Horner's rule and of summation [Hi, (5.3)] added to the enclosure.
\item The box end is the double $\xc=\mathrm{fl}(0.89)=0.88999999999999990\ldots$; the certificate is a statement about the chain with this value of $\xc$.
\end{enumerate}
Under (A1)--(A5) every number reported as an upper bound encloses the corresponding exact quantity of the chain with $\xc=\mathrm{fl}(0.89)$. The implementation was tested against the enclosure property on random inputs (no violation in $4\cdot10^4$ cases) and, as negative controls, versions with one component deliberately broken (the widening removed, the remainder of $\exp$ or of $\mathrm{erfcx}$ dropped, the summation bound removed) were seen to violate it. The enclosures are not a proof inside a proof assistant: the interval evaluations are outside Lean, and what is in Lean is the algebra that turns the table of intervals into the statement about the window (\S\ref{sec:lean}).

\paragraph{Intervals of $\hbar$: the centred form.} Let $I=[\hbar_1,\hbar_2]$, $\delta=\hbar_2/\hbar_1-1$, $\lambda_k=1/(\varepsilon\hbar_k)$ and $\lambda_0=\sqrt{\lambda_1\lambda_2}$, so that $|\lambda-\lambda_0|\le(\delta/2)\lambda_0$ for $\lambda\in[\lambda_2,\lambda_1]$. The slope of the log-density of one step $P_i(u,\cdot)$ is $\lambda G_\lambda(w)$ with $G_\lambda(w)=u-2w+M_\lambda(w)-\varepsilon^2V'(w)$, where $M_\lambda$ is the mean of the next site. Writing
\[
\lambda G=\lambda_0(u-2w)+(\lambda-\lambda_0)\,G+\lambda_0\big(M-\varepsilon^2V'\big)
\]
the first term is the reference Gaussian, integrated exactly by the error function as before, and the width in $\hbar$ enters only through $(\lambda-\lambda_0)G$, whose size on a cell is at most $(\delta/2)\lambda_0\max|G|$: the relative size $\delta/2$ of the true slope, with no power of $\lambda$. The corners of the product of the two intervals give the slope enclosure $s_j^\pm$ on each cell, and Lemma~\ref{lem:M}(c) applies unchanged (Lean: \code{centered\_split}, \code{geomean\_dev}, \code{mul\_le\_corners}). A cell of $b$ is handled by two chains, with $b$ at its lower and upper end (Lemma~\ref{lem:M}(b)), and within each chain only $\lambda\in I$ is enclosed. The recursion is run with $\lambda_1$ for the upper bound and $\lambda_2$ for the lower bound (the lower bound is a check only).

\paragraph{The table.} The range $[10^{-5},1]$ is covered by $239$ intervals, fixed in advance (no adaptation during the run), with relative widths $\delta=\hbar_2/\hbar_1-1$ from $5\cdot10^{-4}$ (next to $\hbar=10^{-2}$) to $0.63$; the covering has no gaps (Lean: \code{geometric\_cover\_count}, \code{cover\_le}, \code{join\_windows}). For each interval the certificate is run for every geometry of the coarse chain (interior $r=7$ with the nine $b$-cells of \S\ref{sec:cells}; ends $r=0,\dots,7$ with $b\in(-\infty,-\xc]$), and the largest upper bound over the geometries is recorded. The run took $33$ CPU-hours on one machine; the record (one entry per interval with its settings, cells, bounds and the worst geometry) was checked independently for completeness, for the absence of gaps, and for identical settings across all $239$ entries.

\subsection{Small \texorpdfstring{$\hbar$}{hbar} on paper}\label{sec:PIN}
For $\hbar\le10^{-5}$ the coarse sites are pressed against the wall $-\xc$ and the influence is controlled on paper.

\begin{proposition}[PIN]\label{prop:PIN}
Let $\lambda=1/(\varepsilon\hbar)$, $L_g:=1980/1217$, $h_{77}:=3664648/1770735$, $D_0:=0.023126$ and $d(\hbar):=D_0-2\sqrt{7h_{77}\varepsilon\hbar}$. For every $0<\hbar\le10^{-5}$ and every geometry of the coarse chain,
\[
0\le\rho(a,b;r)\le\frac{2L_g\,\varepsilon\hbar}{d(\hbar)^2}+\text{(tail)}\le0.437,\qquad\text{hence}\quad\rcol(\hbar)\le0.874<1 .
\]
The bound is increasing in $\hbar$ and $\rho=O(\hbar)$ as $\hbar\to0$. The value $D_0$ is a lower bound, checked by one-dimensional interval arithmetic, of the classical force $D(u)=2u-m_7(u)-m_9(u)+\varepsilon^2V'(u)$ on $[-\xc,-0.8]$ in the worst geometry; everything else is on paper, with Brascamp--Lieb for linear functionals as the cited ingredient; the comparison principle for the classical minimisers is Lemma~\ref{lem:CMP} (Appendix~\ref{app:CMP}).
\end{proposition}
\begin{proof}[Proof (outline)]
\emph{(L0) Convexity of short bridges.} $\varepsilon^2V''\ge-1/8$, so the Hessian of a bridge of $n$ sites with fixed ends is at least $H_0^{(n)}=\mathrm{tridiag}(-1,15/8,-1)$, whose leading principal minors $D_n$ are positive for $n\le7$ and negative for $n=8$ (Lean: \code{D\_pos\_le\_seven}, \code{D\_eight\_neg}; $D_7=1770735/2097152$). Bridges of at most $7$ sites are uniformly convex; the coarse spacing $8$ confines the non-convex region to the inside of a bridge.
\emph{(L1)} Given $z_8=u$ the two halves are independent bridges, $a$ enters only through $e^{-\lambda(z_1-a)^2/2}$, so $\rho=\lambda\Cov(z_8,g(z_8))$ with $g(u)=\E[z_1\mid z_8=u]$ increasing (TP$_2$), hence $\rho\ge0$.
\emph{(L2) Brascamp--Lieb on the bridge.} $g'(u)=\lambda\Cov(z_1,z_7)\le\lambda\Var(z_1+z_7)/4\le L_g$ with $L_g=(h_{11}+h_{77}+2h_{17})/4$, $h_{ij}=((H_0^{(7)})^{-1})_{ij}$ (rational values in Lean: \code{h11\_value}, \code{h17\_value}, \code{Lg\_value}). Thus $\rho\le\lambda L_g\Var(z_8)\le\lambda L_g\,\E[(z_8+\xc)^2]$.
\emph{(L3) Integration by parts on the box.} For $\mu\propto e^{-\lambda U}$ on a convex box with minimiser $x^*$, the divergence theorem (faces by Fubini, corners of measure zero, boundary term $\ge0$ by convexity) gives $\E\langle\nabla U,X-x^*\rangle\le n/\lambda$; with $\mathrm{Hess}\,U\ge H_0$ this yields $\E[z_7]\le m_7+\sqrt{7h_{77}/\lambda}$ for the classical minimiser $m_7$ of the left bridge, and likewise on the right.
\emph{(L4), (L5)} The marginal density of $z_8$ is $\propto e^{-\lambda W(u)}$ with $W'(u)=2u-\E_L[z_7]-\E_R[z_9]+\varepsilon^2V'(u)$. By the comparison principle of Lemma~\ref{lem:CMP} (the minimiser of a bridge is nondecreasing in its boundary values), $m_7$ and $m_9$ are largest in the worst geometry ($a=0$; $b=0$ in the interior, $b=-\xc$ at the ends), so $W'\ge D(u)-2e$ with $e=\sqrt{7h_{77}\varepsilon\hbar}$. On $[-\xc,-0.8]$, $D\ge D_0$ and the conditional law is stochastically dominated by an exponential of rate $\lambda d$, so $\E[(z_8+\xc)^2]\le2/(\lambda d)^2$ there; on $[-0.8,0]$ the mass is at most $\xc^3\lambda M e^{-\lambda\Delta}$ with $\Delta\ge0.0041$ at $\hbar=10^{-5}$ and $\lambda\Delta\ge819$, numerically zero. Assembling (Lean: \code{assemble}, \code{final\_bound}): $\rho\le2L_g\varepsilon\hbar/d^2+\text{tail}\le0.437$ at $\hbar=10^{-5}$ ($e=0.00851$, $d=0.00610$), and the bound increases with $\hbar$.
\end{proof}
The loss against the measured values (about $1.4\cdot10^{-4}$ at $\hbar=10^{-5}$) is three orders of magnitude; it comes from the $\sqrt n$ of (L3) and from not using the pinning of $z_7$. At $\hbar=7\cdot10^{-6}$ the bound gives $\rcol\le0.29$, at $5\cdot10^{-6}$, $0.13$.

\subsection{Proposition DOB}

\begin{proposition}[DOB]\label{prop:DOB}
\textup{(a) [interval arithmetic under (A1)--(A5)]} For every $\hbar\in[10^{-5},1]$ and every geometry of the coarse chain \textup{(}interior $r=7$, $b\in[-\xc,0]$; ends $r=0,\dots,7$, $b\in(-\infty,-\xc]$\textup{)} and every $a\in[-\xc,0]$, $\rho(a,b;r)\le0.1997$; hence $\rcol(\hbar)\le0.400$. The $239$ intervals all pass with largest upper bound $0.199666$, attained in the interior geometry on $[1.325\cdot10^{-4},1.433\cdot10^{-4}]$.

\textup{(b) [paper, plus one one-dimensional interval-arithmetic lower bound]} For every $\hbar\in(0,10^{-5}]$, $\rcol(\hbar)\le0.874$ \textup{(}Proposition~\ref{prop:PIN}\textup{)}.

\textup{(c) [paper, plus one floating-point value; not needed below]} Under the convention that for $k\equiv0\pmod 8$ the coarse sites are $y_m=z_{8m-1}$ \textup{(}so that both neighbours of every coarse site lie in the box\textup{)}, for $\hbar\in[0.27,1]$ one has $\sup\rho\le(V_{\mathrm{in}}(0.27)/(0.27\,\varepsilon))^2$ and hence $\rcol\le0.858$, provided $V_{\mathrm{in}}(0.27)\le0.0884$. Here $V_{\mathrm{in}}(\hbar)=\min_\theta\max(\E_{\nu_R}(X-\theta)^2,\E_{\nu_L}(X-\theta)^2)$ with $\nu_R\propto e^{-(x^2/\varepsilon+\varepsilon V(x))/\hbar}$ on $[\theta,0]$ and $\nu_L\propto e^{-((x+\xc)^2/\varepsilon+\varepsilon V(x))/\hbar}$ on $[-\xc,\theta]$. The value $V_{\mathrm{in}}(0.27)=0.08744$ was computed in floating point by the trapezoidal rule, without interval arithmetic.

In particular $\rcol(\hbar)<1$ for every $\hbar\in(0,1]$.
\end{proposition}

Part (a) is the output of the certificate of \S\ref{sec:discipline}; its proof is the computation, whose inputs, outputs and hashes are distributed (\S\ref{sec:lean}). Table~\ref{tab:cert} summarises the record by decades of $\hbar$. The lower bounds (the same inequality reversed) are reported as a consistency check; in every interval lower $\le$ upper. The interval bounds are coarser than the point bounds of the previous version, which is the price of the width $\delta$; they are not comparable cell by cell, since the cells differ. The method of the previous version (double precision with outward margins) was also run on the same $239$ intervals with the same cells: for $\hbar\ge10^{-4}$ the two upper bounds agree to within $4\cdot10^{-4}$ relatively; below $10^{-4}$ the previous method is larger by up to $20\%$ and, in seven intervals near $3\cdot10^{-5}$ where its worst geometry is an end geometry, by up to a factor $3.4$; in no interval is the interval-arithmetic bound the larger one. Those values are not used.

\begin{proof}[Proof of \textup{(c)}]
For a positive density $p$ on $[l,r]$ with $(\log p)'\le c_R$ on $(\theta,r)$ and $(\log p)'\ge-c_L$ on $(l,\theta)$, $\E_p(X-\theta)^2\le\max(\E_{\nu_R}(X-\theta)^2,\E_{\nu_L}(X-\theta)^2)$, by the Chebyshev order inequality on each side (the ratio $p/\nu$ decreases in $|X-\theta|$ while $(X-\theta)^2$ increases). For the middle site, $-\partial_sS\le-2s/\varepsilon-\varepsilon V'(s)$ (both neighbours at $0$) and $-\partial_sS\ge-[2(s+\xc)/\varepsilon+\varepsilon V'(s)]$ (both at the wall), which are the densities $\nu_R,\nu_L$; the same bounds hold for $z_1$ and $z_7$ in the left block. With the Hoeffding representation $\rho=(\varepsilon\hbar)^{-1}\int\varphi'(t)H(t)\,dt$, $\varphi(s)=\E[z_1\mid z_8=s]$, $H\ge0$, $\int H=\Var_m$ [Hoe], and $\varphi'=\Cov(z_1,z_7\mid s)/(\varepsilon\hbar)\le V_{\mathrm{in}}/(\varepsilon\hbar)$ by Cauchy--Schwarz, one gets $\sup\rho\le(V_{\mathrm{in}}(\hbar)/(\varepsilon\hbar))^2$. The comparison densities are monotone in $\hbar$ in the likelihood-ratio order, so $V_{\mathrm{in}}(\hbar)/(\varepsilon\hbar)$ is non-increasing and the value at $\hbar=0.27$ suffices; $2(0.0884/0.135)^2<0.86$ is checked in Lean. The use of the convention for $k\equiv0\pmod8$ in Theorem~\ref{thm:CC} only needs block lengths $\le7$ and at most $14$ sites around a coarse site, which the shifted division keeps.
\end{proof}

\begin{table}[ht]
\centering\small
\setlength{\tabcolsep}{4pt}\begin{tabular}{@{}lrlrlllr@{}}
\toprule
$\hbar$ & intervals & worst interval & $u$-cells & upper & geometry & lower & CPU h\\
\midrule
$[10^{-5},10^{-4}]$ & 44 & $[8.846\cdot10^{-5},\,9.313\cdot10^{-5}]$ & 2047 & $0.1989$ & interior & $1.3\cdot10^{-5}$ & 20.7\\
$[10^{-4},10^{-3}]$ & 38 & $[1.325\cdot10^{-4},\,1.433\cdot10^{-4}]$ & 1577 & $0.1997$ & interior & $3.3\cdot10^{-5}$ & 6.1\\
$[10^{-3},10^{-2}]$ & 52 & $[1.201\cdot10^{-3},\,1.282\cdot10^{-3}]$ & 996 & $0.1923$ & interior & $6.2\cdot10^{-4}$ & 3.8\\
$[10^{-2},10^{-1}]$ & 95 & $[0.01082,\,0.01125]$ & 420 & $0.1828$ & interior & $3.2\cdot10^{-3}$ & 2.5\\
$[10^{-1},1]$ & 10 & $[0.1,\,0.11]$ & 133 & $0.1361$ & interior & $5.9\cdot10^{-4}$ & $<0.1$\\
\bottomrule
\end{tabular}
\caption{The certificate of Proposition~\ref{prop:DOB}(a), by decades of $\hbar$: number of intervals, the interval with the largest upper bound for $\sup\rho$, the largest number of $u$-cells in the decade, that upper bound, the geometry attaining it, the corresponding lower bound, and the CPU time. Every one of the $239$ intervals has upper bound $\le0.1997$; in $232$ of them the worst geometry is the interior one, in $7$ an end geometry ($r=0$ or $r=3$). The cells near the wall at small $\hbar$ dominate the cost (\S\ref{sec:cells}).}
\label{tab:cert}
\end{table}

\paragraph{Controls.} With the division every fourth site (blocks of length $3$) the point certificate of the previous version returned upper bounds $2.20$, $10.6$, $365$ at $\hbar=0.03,0.01,10^{-3}$, so it does not certify $\rcol<1$ there (necessarily, being an upper bound; this is not a test of power). For the Gaussian chain ($V\equiv0$), whose exact value is $1/2$, the interval-table certificate in double precision returns an enclosure containing $1/2$, which is loose. The negative controls of the interval-arithmetic implementation are described in \S\ref{sec:discipline}.

\begin{corollary}\label{cor:points}
Under the hypotheses of Theorem~\ref{thm:CC} other than \textup{(DOB)}, for all $k\ge1$: for $\hbar\in[10^{-5},1]$, with the qualification of Proposition~\ref{prop:DOB}\textup{(a)}, $C_P(Q_k)\le9.18+e^{87.43}\hbar^{-8}$; for $\hbar\in(0,10^{-5}]$, with Proposition~\ref{prop:DOB}\textup{(b)}, $C_P(Q_k)\le9.18+e^{88.99}\hbar^{-8}$. In particular $C_P(Q_k)\le9.18+e^{88.99}\hbar^{-8}$ for every $\hbar\in(0,1]$.
\end{corollary}
\begin{proof}
$\ln(c_6e^{80.79})=86.9121$; $1/(1-0.400)=e^{0.5109}$ and $1/(1-0.874)=e^{2.0715}$; the exponents are rounded up.
\end{proof}

\subsection{Geometries and cells}\label{sec:cells}
The supremum defining $\rcol$ runs over every conditional law that occurs in the coarse chain. We cover them as follows. Interior coarse sites ($r=7$): $b\in[-\xc,0]$ is split into the nine cells with endpoints $-0.89,-0.6,-0.4,-0.25,-0.15,-0.08,-0.04,-0.02,-0.01,0$. Coarse sites next to the end of the chain ($r=0,\dots,7$): $b=\beta\in(-\infty,-\xc]$ is split into cells whose lowest member is handled by the limit $b\to-\infty$ of Lemma~\ref{lem:M}(b); for $r\le2$ the cells near $-\xc$ have width of order $\sqrt{\varepsilon\hbar}$. The first coarse site is the mirror image of the last. The variable $a\in[-\xc,0]$ is covered by the $u$-cells of width about $c_1\sqrt{\varepsilon\hbar}$, with $c_1\in[0.015,0.25]$ depending on $\hbar$; the column ``$u$-cells'' of Table~\ref{tab:cert} is the largest number used in each decade.

At $\hbar=10^{-5}$ the end geometries with $r\le2$ need narrow $b$-cells: for $r=0$ the centre $(u+b)/2$ of the last kernel moves by half the width of the $b$-cell, and a cell much wider than $\sqrt{\varepsilon\hbar}$ makes the sandwich of Lemma~\ref{lem:M}(c) much wider than the law itself. With widths $0.01$ ($r=0$) and $0.02$ ($r=1$), $50$ cells in total for $r\le2$, these geometries pass; the decade $[10^{-5},10^{-4}]$ accounts for $21$ of the $33$ CPU-hours of the run. This affects only the cost of the computation, not its validity: Lemma~\ref{lem:M}(c) holds for every partition.

\subsection{Where the result sits in the conditional argument}
In the conditional argument for the double-well path measure, the Poincar\'e constant $C_{\mathrm{core}}(\hbar):=\sup_{k,\alpha,\beta}C_P(Q_k)$ enters a second-order term of the form $C_{\mathrm{core}}\cdot D/\xi^2$, where $\xi$ is the tunnelling correlation length, which grows like $e^{c/\hbar}$. Any bound polynomial in $\hbar$ therefore makes that term vanish as $\hbar\to0$, however crude its constant. Theorem~\ref{thm:CC} with Proposition~\ref{prop:DOB} provides such a bound for every $\hbar\in(0,1]$, under the qualifications stated there. The first step of this argument for the $0+1$-dimensional double well rests on seven provisional inputs, of which one is the core Poincar\'e bound and another is a cross-term bound, which does not involve $C_{\mathrm{core}}$. The present result bears on the core Poincar\'e bound only. Of the two ways in which that bound is supported there, the box-side core and the compartment-side core, it bears on the second, whose core has both outer values in the well; the other six inputs, including the cross-term bound, are unchanged.

\section{What is machine-checked}\label{sec:lean}

The following Lean~4 modules (Mathlib [M]) contain the algebraic and finite steps; measure, integration, limits and the analytic lemmas are outside Lean.
\begingroup\raggedright
\begin{itemize}[leftmargin=*]
\item \code{LatticeGauge.HiddenConvexity}: \code{geom\_mean\_midpoint}, \code{act\_eq\_actRho}, \code{actRho\_midpoint}, \code{act\_min\_unique} (K$'$ for $a,b\le0$, midpoint form).
\item \code{LatticeGauge.DoeblinPoincare}: \code{path\_coupling\_step}, \code{poincare\_of\_block\_contraction}, \code{spectral\_weight\_zero}, \code{contraction\_rate\_pos} (Theorem~\ref{thm:B}).
\item \code{LatticeGauge.BlockPoincare}: \code{kappa\_seven\_gt} ($\kappa_0>1/20$), \code{free\_energy\_quasiconvex}, \code{quasiconvex\_add\_osc}, \code{ratio\_exp\_bound}, \code{hardy\_discrete}, \code{two\_scale\_assembly} (Lemmas~\ref{lem:split}--\ref{lem:QC}).
\item \code{LatticeGauge.CoarseChainPoincare}: \code{cov\_le\_lip\_var}, \code{cross\_term}, \code{sq\_three}, \code{coarse\_grad\_sum}, \code{path\_coupling\_site}, \code{poincare\_of\_site\_contraction}, \code{two\_level\_assembly}, \code{coarse\_chain\_poincare}, \code{g7\_value}, \code{assembly\_const} (Theorems~\ref{thm:C}, \ref{thm:CC}).
\item \code{LatticeGauge.InfluenceRecursion}: \code{tp2\_comp}, \code{cdf\_upper}, \code{cdf\_lower}, \code{k\_upper}, \code{step\_bound}, \code{gauss\_product} (Lemmas~\ref{lem:R}, \ref{lem:M}).
\item \code{LatticeGauge.CenterCovariance}: \code{hoeffding\_kernel}, \code{kernel\_le\_sup}, \code{chebyshev\_double}, \code{large\_hbar\_column} (Proposition~\ref{prop:DOB}(c)).
\item \code{LatticeGauge.DobrushinIntervalTable}: \code{centered\_split}, \code{le\_of\_sq\_le\_sq\_nonneg}, \code{geomean\_dev}, \code{mul\_le\_corners}, \code{steps\_growth}, \code{one\_add\_pow\_le\_exp}, \code{geometric\_cover\_count}, \code{cover\_le}, \code{window\_le\_of\_interval\_table}, \code{column\_lt\_one\_of\_table}, \code{join\_windows} (the centred form and the passage from the table of intervals to the window, \S\ref{sec:discipline}).
\item \code{LatticeGauge.PinnedCoarseChain}: \code{D\_seven}, \code{D\_six}, \code{D\_pos\_le\_seven}, \code{D\_eight\_neg}, \code{h11\_value}, \code{h17\_value}, \code{Lg\_value}, \code{assemble}, \code{final\_bound} (the rational constants and the assembly of Proposition~\ref{prop:PIN}).
\end{itemize}
Each statement carries its hypotheses in the Lean signature. The independent re-check of the modules was completed on 2026-10-15, and that of the two modules added in this version on 2026-10-17 (every file checked with \code{lake env lean} in the same environment on a separate session, rc=0, standard axioms only). The bundle is listed at \url{\bundleurl} (with a \code{.sha256} file at the same address) and the complete \code{\#print axioms} output at \url{\axiomsurl}; both addresses were placed on 2026-10-15 and updated on 2026-10-17, after the independent re-check. To check: unpack, run \code{lake exe cache get}, then \code{LEAN\_NUM\_THREADS=1 lake build} on each module above, and compare the trailing \code{\#print axioms} lines with the log (expected: \code{propext}, \code{Classical.choice}, \code{Quot.sound} only). The ledger and build instructions of the site's Lean developments are at \url{https://computoergosum.com/en/principia/lean/index.html}. The certificate's programs, specification files and output files (one JSON per run, with every cell's upper and lower value), a manifest of the files behind Table~\ref{tab:cert}, and SHA-256 hashes are part of the same bundle. In the repository named in \S\ref{sec:ai} these files are kept under \code{papers/\allowbreak{}ym-coarse-chain-poincare/\allowbreak{}certificates} (\code{certificates/\allowbreak{}interval/}, with the table of $239$ intervals, the merged output of the four processes and their logs; the files of version~1 remain in \code{certificates/}; the one-dimensional interval computation of Proposition~\ref{prop:PIN} in \code{certificates/\allowbreak{}pin/}).\par\endgroup

\section{Correspondence table}\label{sec:corr}

Line numbers refer to the files \code{LatticeGauge<Module>.lean} of the bundle; the third column gives \code{<Module>}: line of each declaration in the order of the second column.

\begin{center}\footnotesize
\begin{tabular}{@{}p{0.20\linewidth}p{0.47\linewidth}p{0.27\linewidth}@{}}
\toprule
paper & Lean declaration (namespace \code{LatticeGauge.}) and what is outside & file: line\\
\midrule
(K$'$), \S\ref{sec:setting} & \code{actRho\_midpoint}, \code{act\_min\_unique}; partial minimisation on paper & \code{HiddenConvexity}: 75, 115\\
Theorem~\ref{thm:B} & \code{path\_coupling\_step}, \code{spectral\_weight\_zero}, \code{poincare\_of\_block\_contraction}; FKG, Strassen, spectral theorem on paper & \code{DoeblinPoincare}: 16, 54, 29\\
Lemma~\ref{lem:BE7} & \code{kappa\_seven\_gt}; Bakry--\'Emery cited & \code{BlockPoincare}: 26\\
Lemma~\ref{lem:split} & \code{two\_scale\_assembly}, \code{cross\_term\_sq}, \code{deriv\_sq\_split}; total variance and differentiation on paper & \code{BlockPoincare}: 101, 88, 96\\
Lemma~\ref{lem:H1D} & \code{hardy\_discrete}, \code{ratio\_exp\_bound}; Fubini on paper & \code{BlockPoincare}: 137, 66\\
Lemma~\ref{lem:QC} & \code{free\_energy\_quasiconvex}, \code{quasiconvex\_add\_osc}; Laplace-type bounds on paper & \code{BlockPoincare}: 37, 51\\
Lemma~\ref{lem:B1}, Lemma~\ref{lem:M}(b) & \code{tp2\_comp} (finite form) & \code{InfluenceRecursion}: 18\\
Theorem~\ref{thm:C} & \code{path\_coupling\_site}, \code{poincare\_of\_site\_contraction} & \code{CoarseChainPoincare}: 92, 98\\
Theorem~\ref{thm:CC} & \code{cross\_term}, \code{sq\_three}, \code{coarse\_grad\_sum}, \code{two\_level\_assembly}, \code{coarse\_chain\_poincare}, \code{g7\_value}, \code{g7\_lt}, \code{assembly\_const}; BL cited & \code{CoarseChainPoincare}: 51, 47, 68, 78, 117, 128, 133, 136\\
Lemma~\ref{lem:R} & \code{cdf\_upper}, \code{cdf\_lower}, \code{k\_upper}, \code{step\_bound}, \code{gauss\_product} & \code{InfluenceRecursion}: 49, 56, 60, 68, 76\\
Proposition~\ref{prop:N1} & paper only & ---\\
Centred form, \S\ref{sec:discipline} & \code{centered\_split}, \code{geomean\_dev}, \code{mul\_le\_corners}, \code{steps\_growth}, \code{one\_add\_pow\_le\_exp} & \code{DobrushinIntervalTable}: 20, 32, 51, 68, 79\\
Proposition~\ref{prop:DOB}(a) & \code{geometric\_cover\_count}, \code{cover\_le}, \code{window\_le\_of\_interval\_table}, \code{column\_lt\_one\_of\_table}, \code{join\_windows}; the interval evaluations outside Lean (interval arithmetic under (A1)--(A5)) & \code{DobrushinIntervalTable}: 86, 97, 111, 118, 127\\
Proposition~\ref{prop:PIN} & \code{D\_seven}, \code{D\_six}, \code{D\_pos\_le\_seven}, \code{D\_eight\_neg}, \code{h11\_value}, \code{h17\_value}, \code{Lg\_value}, \code{assemble}, \code{final\_bound}; Brascamp--Lieb, integration by parts and the comparison principle on paper; $D_0$ outside (one-dimensional interval arithmetic) & \code{PinnedCoarseChain}: 17, 20, 24, 29, 33, 36, 40, 45, 51\\
Proposition~\ref{prop:DOB}(c) & \code{hoeffding\_kernel}, \code{kernel\_le\_sup}, \code{chebyshev\_double}, \code{large\_hbar\_column}; the value $V_{\mathrm{in}}(0.27)$ outside & \code{CenterCovariance}: 44, 143, 186, 206\\
\bottomrule
\end{tabular}
\end{center}

\section{What is not claimed}\label{sec:notclaimed}

\begin{table}[ht]
\centering\small
\begin{tabular}{@{}p{0.13\linewidth}p{0.47\linewidth}p{0.32\linewidth}@{}}
\toprule
row & content & status\\
\midrule
END & outer boundary values $\alpha,\beta\le-\xc$ & hypothesis of every statement\\
K$'$ & convexity of $\tilde S$ in $\rho=z^2$, jointly with $(a^2,b^2)$ & paper (Appendix~\ref{app:K}); algebraic core in Lean\\
BE$_7$ & $C_P\le\varepsilon/\kappa(n)$ for blocks of length $\le7$ & paper ($\kappa_0>1/20$ in Lean) + cited [BoL, Prop.~3.1] (open convex domain, no boundary regularity required), [BGL]\\
BL & $\Var\langle e,z\rangle\le\hbar\langle e,M^{-1}e\rangle$ on convex sets & cited [BL, Thm.~4.1]; for the linear functionals used here it also follows from Pr\'ekopa--Leindler [BoL, (1.1)] in a few lines; $g_7$ in Lean\\
TP$_2$ & monotonicity of the coarse conditional laws & paper [Ka]; finite form in Lean\\
FKG & monotone coupling for blocks (Theorem~\ref{thm:B} only) & standard, cited [Ho], [Pr]\\
H1D, QC & Hardy-type bound; quasi-convexity plus oscillation & paper; discrete cores in Lean\\
OLL & contraction implies gap & cited, checked (arXiv version) [Ol]\\
ALG & coupling step, comparison, assembly, $c_6<497$ & Lean\\
DOB, $[10^{-5},1]$ & $\rcol\le0.400$ & interval arithmetic over $239$ intervals under (A1)--(A5); table-to-window algebra in Lean\\
DOB, $(0,10^{-5}]$ & $\rcol\le0.874$ & paper (Proposition~\ref{prop:PIN}) + one one-dimensional interval-arithmetic lower bound; rational constants and assembly in Lean\\
DOB, $[0.27,1]$ & $\rcol\le0.858$ & paper + one floating-point value (not needed for Corollary~\ref{cor:points})\\
\bottomrule
\end{tabular}
\caption{Status of the rows on which Theorem~\ref{thm:CC} and Corollary~\ref{cor:points} rest.}
\label{tab:status}
\end{table}

\begin{enumerate}[leftmargin=*]
\item \textbf{Assumptions of the certificate.} Proposition~\ref{prop:DOB}(a) rests on (A1)--(A5) of \S\ref{sec:discipline}: the outward widening of the five IEEE operations under round-to-nearest, the interval type of \code{mpmath} for the constants, the remainder bounds for $\mathrm{erfcx}$ and $\mathrm{erf}$ [DLMF, \S7.12(i)], the a priori rounding bounds for Horner's rule and summation [Hi, (5.3)], and $\xc=\mathrm{fl}(0.89)$. The interval evaluations are outside Lean; a proof-assistant verification of the computation is not claimed.
\item \textbf{$\xc=\mathrm{fl}(0.89)$.} The certificate concerns the chain whose box end is the double nearest to $0.89$; the difference to $0.89$ itself is not treated.
\item \textbf{Small $\hbar$.} On $(0,10^{-5}]$ the bound $\rcol\le0.874$ of Proposition~\ref{prop:PIN} is three orders of magnitude above the measured influence; the loss is in the $\sqrt n$ of the integration-by-parts step. The constant $e^{88.99}$ of Corollary~\ref{cor:points} comes from this part; on $[10^{-5},1]$ it is $e^{87.43}$.
\item \textbf{Cited statements.} BE$_7$ (Bakry--\'Emery on convex domains, [BoL, Prop.~3.1]), BL (Brascamp--Lieb on convex domains) and the FKG/Holley step of Theorem~\ref{thm:B} are standard results used as cited; they are not formalised in Lean. The comparison principle for the classical minimisers used in Proposition~\ref{prop:PIN} is proved in Appendix~\ref{app:CMP} (Lemma~\ref{lem:CMP}); it is not cited.
\item \textbf{Constants.} $e^{87}\hbar^{-8}$ is crude; measured block and site constants are of order $10^2$ and below $0.27$ respectively. Only a Poincar\'e inequality is proved, not a log-Sobolev inequality.
\item \textbf{One dimension.} The model is a one-dimensional chain (a $0+1$-dimensional path measure); nothing is claimed about four-dimensional lattice gauge theory. In the conditional argument for this double well, whose first step keeps seven provisional inputs, this paper bears only on the core Poincar\'e bound, through its compartment-side support, and leaves the other six, including the cross-term bound, unchanged.
\item \textbf{No external review.} The proofs, the computations and this text were produced within the authors' group (\S\ref{sec:ai}).
\end{enumerate}

\section{Disclosure of AI use}\label{sec:ai}

This note was produced with the assistance of AI agents. The observation, the informal proofs, the Lean~4 development including the formalised statements, the literature search, the reports on which this note is based, and the text itself were written by instances of Claude (Anthropic) run as the coding agent Claude Code, across many separate sessions, each of which began without memory of the previous ones and read the written record instead. A second group of instances, running on the server of the website named on the title page, is responsible for rebuilding the Lean development from the distributed archive on a separate machine, for checking the claims of this note against the reports, and for preparing the distributed version; the independent rebuild was completed on 2026-10-15 (every file checked with \code{lake env lean} in the same environment on a separate session, rc=0, standard axioms only); for version~2 the same group re-checked the two added modules in the same way on 2026-10-17 and verified the record of the interval run (239 entries, no gaps, all passed, identical settings) against the state file of the computation. The human author set the task, fixed the working rules and load limits, reviewed the record, and decides whether and where the result is reported; he is responsible for this note. Adversarial review was carried out only by further instances of the same model family, which share the same blind spots. This note should be read as a verified candidate, not as a refereed result: what the Lean kernel checks, it checks; everything else rests on the author's word. In particular the analytic steps listed in \S\ref{sec:lean} are checked on paper only. This edition and the Lean development it cites are kept in the repository \url{https://github.com/yoshikazuhosomura-sys/principia} (the note under \texttt{papers/\allowbreak{}ym-coarse-chain-poincare}, the development under \texttt{lean/\allowbreak{}ym-coarse-chain-poincare}). The web edition of this note, published on the site named on the title page, lists the same work under the names Kiichi, Shiori and Rin, the human author and the two groups of AI agents; the two editions differ only in the title page, in this section, in the naming of the agent that carried out the independent re-check, and in one sentence of \S\ref{sec:lean} giving the location of the certificate files in the repository.

Version~2 (2026-10-17) adds the interval-arithmetic certificate over $[10^{-5},1]$ (Proposition~\ref{prop:DOB}(b)) and Proposition~\ref{prop:PIN}. The certificate was run on one machine ($33$ CPU-hours, four processes on disjoint sets of intervals, merged); its merged output, the Lean modules of version~2 and the arithmetic of the constants were checked by a second instance of the agent, independent of the one that produced them.

\appendix
\section{Hidden convexity}\label{app:K}

\begin{lemma}\label{lem:K}
Let $n\ge1$, $\varepsilon>0$, $\xc>0$, put $z_0=a$, $z_{n+1}=b$ with $a,b\le0$, and let $S(z)=\sum_{i=0}^{n}(z_{i+1}-z_i)^2/(2\varepsilon)+\varepsilon\sum_{i=1}^{n}V(z_i)$ with $V(z)=(z^2-1)^2/8$, on the box $[-\xc,0]^n$. Write $\rho_i=z_i^2$ ($0\le i\le n+1$) and $\tilde S(\rho):=S(-\sqrt\rho)$.
\begin{enumerate}[label=(\alph*)]
\item $(\rho_0,\rho_1,\dots,\rho_n,\rho_{n+1})\mapsto\tilde S$ is convex on $[0,\infty)^{n+2}$.
\item For fixed $a,b\le0$, $\tilde S$ is $(\varepsilon/4)$-strongly convex in $(\rho_1,\dots,\rho_n)\in[0,\xc^2]^n$.
\item For fixed $a,b\le0$, $S$ has exactly one local minimum on $[-\xc,0]^n$, and it is the global minimum.
\end{enumerate}
\end{lemma}

\begin{proof}
For $z_i,z_{i+1}\le0$, $(z_{i+1}-z_i)^2=\rho_{i+1}+\rho_i-2\sqrt{\rho_i\rho_{i+1}}$. The geometric mean $(x,y)\mapsto\sqrt{xy}$ is concave on $[0,\infty)^2$, so $-\sqrt{\rho_i\rho_{i+1}}$ is convex; this includes the end terms $i=0$ and $i=n$, in which $\rho_0=a^2$ and $\rho_{n+1}=b^2$ enter as variables. The terms linear in $\rho$ are convex, and $\varepsilon V(z_i)=\varepsilon(\rho_i-1)^2/8$ has second derivative $\varepsilon/4$ in $\rho_i$. Since $1/(2\varepsilon)>0$, $\tilde S$ is a sum of convex functions, which gives (a), and of $\varepsilon\sum_{i=1}^n(\rho_i-1)^2/8$, which gives (b). For (c), the componentwise map $\rho\mapsto-\sqrt\rho$ is a homeomorphism of $[0,\xc^2]^n$ onto $[-\xc,0]^n$, so local minima of $S$ and of $\tilde S$ correspond. A local minimum of a strictly convex function on a convex set is the global minimum and there is at most one; $S$ is continuous on the compact box, so the minimum exists.
\end{proof}


\section{Comparison of minimisers}\label{app:CMP}

\begin{lemma}\label{lem:CMP}
Let $n$, $\varepsilon$, $\xc$, $V$ and $S$ be as in Lemma~\ref{lem:K}, and write $S_{a,b}$ for $S$ with boundary values $z_0=a$, $z_{n+1}=b$. Let $a\le a'\le0$ and $b\le b'\le0$, and let $m$ and $m'$ be the minimisers of $S_{a,b}$ and $S_{a',b'}$ on $[-\xc,0]^n$ \textup{(}unique by Lemma~\ref{lem:K}(c)\textup{)}. Then $m_i\le m'_i$ for every $i$.
\end{lemma}

\begin{proof}
Write $x\wedge y$ and $x\vee y$ for the componentwise minimum and maximum; both lie in the box when $x$ and $y$ do. For a convex $h:\R\to\R$, the function $(u,v)\mapsto h(u-v)$ satisfies
\begin{equation}\label{eq:submod}
h(u\wedge u'-v\wedge v')+h(u\vee u'-v\vee v')\le h(u-v)+h(u'-v')\qquad(u,v,u',v'\in\R).
\end{equation}
If $(u,v)$ and $(u',v')$ are ordered this is an equality. Otherwise, say $u\le u'$ and $v\ge v'$; the left-hand arguments $u-v'$ and $u'-v$ lie in the interval $[u-v,\,u'-v']$ and have the same sum as its end points, so \eqref{eq:submod} is the convexity of $h$. Now put $h(t)=t^2/(2\varepsilon)$ and compare $S_{a,b}(m\wedge m')+S_{a',b'}(m\vee m')$ with $S_{a,b}(m)+S_{a',b'}(m')$ term by term. The potential terms agree, since $\{m_i\wedge m'_i,\,m_i\vee m'_i\}=\{m_i,m'_i\}$. Each inner pair term $h(z_{i+1}-z_i)$, $1\le i<n$, obeys \eqref{eq:submod}. The end terms are $h(z_1-a)$ and $h(b-z_n)$; since $a\le a'$ and $b\le b'$, the pairs $(m_1,a)$, $(m'_1,a')$ have meet $(m_1\wedge m'_1,a)$ and join $(m_1\vee m'_1,a')$, and likewise at the other end, so \eqref{eq:submod} applies to them too. Hence
\[
S_{a,b}(m\wedge m')+S_{a',b'}(m\vee m')\le S_{a,b}(m)+S_{a',b'}(m').
\]
As $m$ minimises $S_{a,b}$, $S_{a,b}(m)\le S_{a,b}(m\wedge m')$, so $S_{a',b'}(m\vee m')\le S_{a',b'}(m')$. Thus $m\vee m'$ is a minimiser of $S_{a',b'}$, and by uniqueness $m\vee m'=m'$, that is, $m\le m'$.
\end{proof}

In Proposition~\ref{prop:PIN} the lemma is applied to the left bridge $z_1,\dots,z_7$ (boundary values $a$ and $z_8=u$) and to the right bridge (boundary values $u$ and $b$; for $r=0$ there is no right bridge and $m_9=b$): for fixed $u$ and $r$, $m_7$ is largest at $a=0$, and $m_9$ at the largest admissible $b$, namely $b=0$ in the interior and $b=-\xc$ at the ends.
\begin{thebibliography}{KSR2}
\bibitem[BGL]{BGL} D.~Bakry, I.~Gentil and M.~Ledoux, \emph{Analysis and Geometry of Markov Diffusion Operators}, Grundlehren Math. Wiss. 348, Springer, 2014, DOI 10.1007/978-3-319-00227-9; cited for the curvature criterion on convex domains (not consulted directly for this statement).
\bibitem[BL]{BL} H.~J.~Brascamp and E.~H.~Lieb, \emph{On extensions of the Brunn--Minkowski and Pr\'ekopa--Leindler theorems, including inequalities for log concave functions, and with an application to the diffusion equation}, J. Funct. Anal. 22 (1976), 366--389, DOI 10.1016/0022-1236(76)90004-5; cited for Thm.~4.1 (not consulted directly).
\bibitem[BoL]{BoL} S.~G.~Bobkov and M.~Ledoux, \emph{From Brunn--Minkowski to Brascamp--Lieb and to logarithmic Sobolev inequalities}, Geom. Funct. Anal. 10 (2000), 1028--1052.
\bibitem[DLMF]{DLMF} NIST Digital Library of Mathematical Functions, \S7.12(i), \url{https://dlmf.nist.gov/7.12}.
\bibitem[Do]{Do} R.~L.~Dobrushin, \emph{Prescribing a system of random variables by conditional distributions}, Theory Probab. Appl. 15 (1970), 458--486, DOI 10.1137/1115049 (not consulted directly).
\bibitem[Hoe]{Hoe} W.~Hoeffding, \emph{Masstabinvariante Korrelationstheorie}, Schr. Math. Inst. Univ. Berlin 5 (1940), 181--233 (not consulted directly).
\bibitem[Hi]{Hi} N.~J.~Higham, \emph{Accuracy and Stability of Numerical Algorithms}, 2nd ed., SIAM, 2002.
\bibitem[Ho]{Ho} R.~Holley, \emph{Remarks on the FKG inequalities}, Comm. Math. Phys. 36 (1974), 227--231, DOI 10.1007/BF01645980 (not consulted directly).
\bibitem[HS]{HS} R.~Holley and D.~Stroock, \emph{Logarithmic Sobolev inequalities and stochastic Ising models}, J. Stat. Phys. 46 (1987), 1159--1194, DOI 10.1007/BF01011161 (not consulted directly).
\bibitem[Ka]{Ka} S.~Karlin, \emph{Total Positivity}, Vol.~I, Stanford University Press, 1968 (not consulted directly).
\bibitem[KSR]{KSR} Kiichi, Shiori and Rin, \emph{Yang--Mills and the mass gap} (online textbook, in Japanese), \url{https://computoergosum.com/principia/yang-mills/index.html} (2026).
\bibitem[KSR2]{KSR2} Kiichi, Shiori and Rin, \emph{The inverse squared spectral gap bounds the Langevin Poincar\'e constant of Euclidean path measures from below}, 2026-10-09, \url{https://computoergosum.com/en/index.html}.
\bibitem[Ma]{Ma} F.~Martinelli, \emph{Lectures on Glauber dynamics for discrete spin models}, in: Lectures on Probability Theory and Statistics (Saint-Flour 1997), Lecture Notes in Math. 1717, Springer, 1999, 93--191, DOI 10.1007/978-3-540-48115-7\_2.
\bibitem[M]{M} Mathlib, \url{https://github.com/leanprover-community/mathlib4}.
\bibitem[mp]{mp} The mpmath development team, \emph{mpmath: a Python library for arbitrary-precision floating-point arithmetic}, \url{https://mpmath.org/}.
\bibitem[Mu]{Mu} B.~Muckenhoupt, \emph{Hardy's inequality with weights}, Studia Math. 44 (1972), 31--38 (not consulted directly).
\bibitem[Ol]{Ol} Y.~Ollivier, \emph{Ricci curvature of Markov chains on metric spaces}, J. Funct. Anal. 256 (2009), 810--864, DOI 10.1016/j.jfa.2008.11.001, arXiv:math/0701886; cited (arXiv v4) for Prop.~29 and Cor.~30.
\bibitem[Pr]{Pr} C.~J.~Preston, \emph{A generalization of the FKG inequalities}, Comm. Math. Phys. 36 (1974), 233--241, DOI 10.1007/BF01645981 (not consulted directly).
\bibitem[SZ]{SZ} D.~W.~Stroock and B.~Zegarlinski, \emph{The logarithmic Sobolev inequality for discrete spin systems on a lattice}, Comm. Math. Phys. 149 (1992), 175--193, DOI 10.1007/BF02096629.
\end{thebibliography}
\end{document}
